%
% menu_lifecycle_delivery_reorder_buggy.tex
%
% FIDELITY CONTROL for docs/menu_lifecycle_delivery.tex.
%
% This spec differs from the fixed menu_lifecycle_delivery.tex in exactly one
% place: the two departure legs run in the WRONG order. The cascade's inbox
% teardown (DepartInboxFirst) runs BEFORE the registry discard
% (DepartRegistryAfter), so the inbox -- and any undrained menu event in it --
% is torn down while the connection is STILL in the registry. A click reported
% delivered is then discarded from a live recipient that still holds a recv():
% a genuine live-recipient loss, M1b's violation.
%
% This is the counterfactual the fixed model's ordering rules out. The shipped
% code guarantees the fixed order -- HubDisplay._depart runs
% self._clients.discard(c) (registry) BEFORE self._cascade.run(c) (which fires
% inbox.drop_session), under one StoreLock hold, and the cascade's docstring
% states "connection_id has already been irreversibly removed from the registry
% by the time this runs". This control removes that guarantee to show it is
% load-bearing: reverse the legs and the loss becomes reachable.
%
% Expected: M1b's negation is reachable.
%   probcli docs/menu_lifecycle_delivery_reorder_buggy.tex -model_check -nodead \
%     -goal "lostLive = fset" -p DEFAULT_SETSIZE 3
%   => goal FOUND
%      (Connect(c); Arm(c); MenuDeliver(c); DepartInboxFirst(c)).
%
% Against the fixed spec the same goal is NOT found.
%

\documentclass[a4paper,10pt,fleqn]{article}
\usepackage[margin=1in]{geometry}
\usepackage{fuzz}
\usepackage[colorlinks=true,linkcolor=blue,citecolor=blue,urlcolor=blue]{hyperref}

\begin{document}

\title{lux Agent-Menu Delivery --- Reorder Fidelity Control}
\author{The cascade that tears down the inbox before the registry discard}
\date{September 2026}
\maketitle

This is a fidelity control for \texttt{menu\_lifecycle\_delivery.tex}. Only the
departure legs differ: they run in the wrong order. Every other schema is
copied verbatim.

\begin{zed}
  [CONN]
\end{zed}

\begin{zed}
  DSTATE ::= dnone | ddelivered | dgone
\end{zed}

\begin{zed}
  FSTATE ::= fclear | fset
\end{zed}

\begin{zed}
  TPHASE ::= tidle | tinbox
\end{zed}

\begin{schema}{MenuFlow}
  registered : \power CONN \\
  hasInbox : \power CONN \\
  menuOwner : \power CONN \\
  pending : \power CONN \\
  reported : DSTATE \\
  lostLive : FSTATE \\
  departPhase : TPHASE \\
  departFor : \power CONN
\where
  \# departFor \leq 1
\end{schema}

\begin{schema}{Init}
  MenuFlow'
\where
  registered' = \emptyset \\
  hasInbox' = \emptyset \\
  menuOwner' = \emptyset \\
  pending' = \emptyset \\
  reported' = dnone \\
  lostLive' = fclear \\
  departPhase' = tidle \\
  departFor' = \emptyset
\end{schema}

\begin{schema}{Connect}
  \Delta MenuFlow \\
  c? : CONN
\where
  c? \notin departFor \\
  registered' = registered \cup \{ c? \} \\
  hasInbox' = hasInbox \\
  menuOwner' = menuOwner \\
  pending' = pending \\
  reported' = reported \\
  lostLive' = lostLive \\
  departPhase' = departPhase \\
  departFor' = departFor
\end{schema}

\begin{schema}{Arm}
  \Delta MenuFlow \\
  c? : CONN
\where
  c? \in registered \\
  hasInbox' = hasInbox \cup \{ c? \} \\
  registered' = registered \\
  menuOwner' = menuOwner \\
  pending' = pending \\
  reported' = reported \\
  lostLive' = lostLive \\
  departPhase' = departPhase \\
  departFor' = departFor
\end{schema}

\begin{schema}{Own}
  \Delta MenuFlow \\
  c? : CONN
\where
  c? \in registered \\
  hasInbox' = hasInbox \cup \{ c? \} \\
  menuOwner' = menuOwner \cup \{ c? \} \\
  registered' = registered \\
  pending' = pending \\
  reported' = reported \\
  lostLive' = lostLive \\
  departPhase' = departPhase \\
  departFor' = departFor
\end{schema}

\begin{schema}{MenuDeliver}
  \Delta MenuFlow \\
  c? : CONN
\where
  c? \in hasInbox \\
  pending' = pending \cup \{ c? \} \\
  reported' = ddelivered \\
  registered' = registered \\
  hasInbox' = hasInbox \\
  menuOwner' = menuOwner \\
  lostLive' = lostLive \\
  departPhase' = departPhase \\
  departFor' = departFor
\end{schema}

\begin{schema}{MenuProviderGone}
  \Delta MenuFlow \\
  c? : CONN
\where
  c? \notin hasInbox \\
  reported' = dgone \\
  registered' = registered \\
  hasInbox' = hasInbox \\
  menuOwner' = menuOwner \\
  pending' = pending \\
  lostLive' = lostLive \\
  departPhase' = departPhase \\
  departFor' = departFor
\end{schema}

\begin{schema}{Drain}
  \Delta MenuFlow \\
  c? : CONN
\where
  c? \in hasInbox \\
  c? \in pending \\
  pending' = pending \setminus \{ c? \} \\
  registered' = registered \\
  hasInbox' = hasInbox \\
  menuOwner' = menuOwner \\
  reported' = reported \\
  lostLive' = lostLive \\
  departPhase' = departPhase \\
  departFor' = departFor
\end{schema}

% -- THE DEFECT: the inbox teardown runs BEFORE the registry discard ---------

% DepartInboxFirst models inbox.drop_session firing as the FIRST departure leg,
% while the connection is still in the registry. If it discards an undrained
% pending event, that event belonged to a still-live recipient -- a genuine
% live-recipient loss, so lostLive becomes fset.
\begin{schema}{DepartInboxFirst}
  \Delta MenuFlow \\
  c? : CONN
\where
  c? \in registered \\
  departPhase = tidle \\
  hasInbox' = hasInbox \setminus \{ c? \} \\
  pending' = pending \setminus \{ c? \} \\
  (c? \in pending \implies lostLive' = fset) \\
  (\lnot (c? \in pending) \implies lostLive' = lostLive) \\
  departPhase' = tinbox \\
  departFor' = \{ c? \} \\
  registered' = registered \\
  menuOwner' = menuOwner \\
  reported' = reported
\end{schema}

% DepartRegistryAfter completes the deregister after the inbox was already gone.
\begin{schema}{DepartRegistryAfter}
  \Delta MenuFlow \\
  c? : CONN
\where
  departPhase = tinbox \\
  c? \in departFor \\
  registered' = registered \setminus \{ c? \} \\
  menuOwner' = menuOwner \setminus \{ c? \} \\
  departPhase' = tidle \\
  departFor' = \emptyset \\
  hasInbox' = hasInbox \\
  pending' = pending \\
  reported' = reported \\
  lostLive' = lostLive
\end{schema}

\end{document}
